\documentclass[a4paper]{arrowhead}

\usepackage[yyyymmdd]{datetime}
\usepackage{etoolbox}
\usepackage[utf8]{inputenc}
\usepackage[T1]{fontenc}
\usepackage{multirow}
\usepackage{longtable}
\usepackage{tikz}
\usepackage{adjustbox}
\AtBeginEnvironment{longtable}{\small}
\BeforeBeginEnvironment{tikzpicture}{\begin{adjustbox}{max width=\textwidth}}
\AfterEndEnvironment{tikzpicture}{\end{adjustbox}}
\usetikzlibrary{positioning,shapes.multipart,arrows.meta}

\renewcommand{\dateseparator}{-}
\setlength{\emergencystretch}{3em}
\lstset{language=HTTP,breakatwhitespace=false,columns=fullflexible}

%% Special references
\newcommand{\mref}[2][]{{\textcolor{ArrowheadPurple}{\hyperref[sec:model:#2]{\ifstrempty{#1}{#2}{#1}}}}}
\newcommand{\oref}[1]{{\textcolor{ArrowheadBlue}{\hyperref[sec:operations:#1]{#1}}}}
\newcommand{\pdef}[1]{{\textcolor{ArrowheadGrey}{#1\label{sec:model:primitives:#1}\label{sec:model:primitives:#1s}\label{sec:model:primitives:#1es}}}}
\newcommand{\pref}[1]{{\textcolor{ArrowheadGrey}{\hyperref[sec:model:primitives:#1]{#1}}}}

% as template_sd.tex; the input and output arguments are passed already formatted (\mref or \pref)
\newrobustcmd\osubsection[3]{
  \addtocounter{subsection}{1}
  \addcontentsline{toc}{subsection}{\protect\numberline{\thesubsection}operation \textcolor{ArrowheadBlue}{#1}}
  \subsection*{
    \thesubsection\quad
    operation
    \textcolor{ArrowheadBlue}{#1}
    (\notblank{#2}{#2}{})
    \notblank{#3}{: #3}{}
  }
  \label{sec:operations:#1}
}
\newrobustcmd\msubsection[2]{
  \addtocounter{subsection}{1}
  \addcontentsline{toc}{subsection}{\protect\numberline{\thesubsection}#1 \textcolor{ArrowheadPurple}{#2}}
  \subsection*{\thesubsection\quad#1 \textcolor{ArrowheadPurple}{#2}}
  \label{sec:model:#2} \label{sec:model:#2es}
}
%%

\usepackage{xurl}
\begin{document}
\sloppy

%% Arrowhead Document Properties
\ArrowheadTitle{authorizationTokenManagement}
\ArrowheadServiceID{authorizationTokenManagement}
\ArrowheadType{Service Description}
\ArrowheadTypeShort{SD}
\ArrowheadVersion{5.3.0}
\ArrowheadDate{2026-11-09}
\ArrowheadAuthor{Jerker Delsing}
\ArrowheadStatus{REVIEW DRAFT}
\ArrowheadContact{jerker.delsing@ltu.se}
\ArrowheadFooter{\href{www.arrowhead.eu}{www.arrowhead.eu}}
\ArrowheadSetup
%%

%% Front Page
\begin{center}
  \vspace*{1cm}
  \huge{\arrowtitle}

  \vspace*{0.2cm}
  \LARGE{\arrowtype}
  \vspace*{1cm}

  \Large{Service ID: \textit{"\arrowid"}}

  \vspace*{\fill}

  % Front Page Image

  \vspace*{1cm}
  \vspace*{\fill}

  % Front Page Abstract
  \begin{abstract}
    This document describes the \texttt{authorization\allowbreak{}Token\allowbreak{}Management} service, version 1.0.0, produced by the ConsumerAuthorization core system of Arrowhead (documentation edition 5.3.0).
  \end{abstract}

  \vspace*{1cm}
\end{center}

\newpage
%%

%% Table of Contents
\tableofcontents
\newpage
%%

\section{Documentation edition and implementation scope}
This is the 5.3.0 documentation edition, dated 2026-11-09. It carries forward the system and service contracts from the retained 5.2.2 documentation and identifies the current primary implementation sources. The former \path|documentation/v5.2.2| folder is retained unchanged.

Document version 5.3.0 is distinct from executable, dependency, registered system/service and wire-contract versions. Values still defined by the source (including 5.2.2 artefact versions, 1.0.0 service versions and \path|SysMonitor-AMQP-IDD-5.2.2| payload identifiers) are preserved. This documentation update does not change runtime behaviour or establish a new implementation release. Historical verification statements describe the tests recorded in the inherited baseline; this revision validates document compilation and references.

\section{Overview}
\label{sec:overview}

This document describes the \texttt{authorization\allowbreak{}Token\allowbreak{}Management} service, which provides operator management of authorization tokens and of provider encryption keys.
It is produced by the ConsumerAuthorization system \cite{ca-sysd} in service version 1.0.0.

The rest of this document is organized as follows.
In Section \ref{sec:operations}, we describe the abstract message operations provided by the service.
In Section \ref{sec:model}, we end the document by presenting the data types used by the mentioned operations.
The concrete interface is specified in the HTTP/TLS/JSON IDD \cite{ca-authorizationTokenManagement-http-idd}. The MQTT binding is specified in the MQTT/TLS/JSON IDD \cite{ca-authorizationTokenManagement-mqtt-idd}.

\subsection{Significant Prior Art}

New in Arrowhead 5.

\subsection{How This Service Is Meant to Be Used}

Dynamic\allowbreak{}Service\allowbreak{}Orchestration calls \texttt{generate-\allowbreak{}tokens} on behalf of consumers during orchestration (it is in \texttt{unbounded.\allowbreak{}token.\allowbreak{}generation.\allowbreak{}whitelist}). Operators query and revoke tokens and register or remove encryption keys of providers.

\subsection{Important Delimitations and Dependencies}
\label{sec:delimitations}

Unbound token generation (bypassing the authorization check) is only allowed for systems in \texttt{unbounded.\allowbreak{}token.\allowbreak{}generation.\allowbreak{}whitelist}.

\newpage

\section{Service Interface}
\label{sec:operations}

This section describes the interfaces to the \texttt{authorization\allowbreak{}Token\allowbreak{}Management} service.
In particular, each subsection names an abstract operation, an input type and an output type, in that order.
The input type is named inside parentheses, while the output type is preceded by a colon.
Input and output types are only denoted when accepted or returned, respectively, by the interface in question.

All abstract data types named in this section are defined in Section \ref{sec:model}.

\begin{figure}[ht!]
  \centering
  \begin{tikzpicture}[font=\small]
    \node[draw=ArrowheadBlue, very thick, rectangle split, rectangle split parts=2, align=center,
          rectangle split part fill={ArrowheadBlue!12,white}, text width=9cm, inner sep=6pt] (b)
      {\guillemotleft{}Service\guillemotright{}\\\textbf{authorizationTokenManagement}
        \nodepart{second} \texttt{generate-\allowbreak{}tokens}\\ \texttt{query-\allowbreak{}tokens}\\ \texttt{revoke-\allowbreak{}tokens}\\ \texttt{add-\allowbreak{}encryption-\allowbreak{}keys}\\ \texttt{remove-\allowbreak{}encryption-\allowbreak{}keys}};
  \end{tikzpicture}
  \caption{SysML-style block of the \texttt{authorization\allowbreak{}Token\allowbreak{}Management} service and its operations.}
  \label{fig:SD}
\end{figure}

The following interface operations are available.

\osubsection{generate-tokens}{\mref[Authorization\allowbreak{}Token\allowbreak{}Generation\allowbreak{}Mgmt\allowbreak{}List\allowbreak{}Request]{AuthorizationTokenGenerationMgmtListRequest}}{\mref[Authorization\allowbreak{}Token\allowbreak{}Mgmt\allowbreak{}List\allowbreak{}Response]{AuthorizationTokenMgmtListResponse}}

Generates tokens on behalf of consumers. With \texttt{unbound=\allowbreak{}true} the authorization check is skipped (only for systems in \texttt{unbounded.\allowbreak{}token.\allowbreak{}generation.\allowbreak{}whitelist}).

Optional parameter: \texttt{unbound}.

\osubsection{query-tokens}{\mref[Authorization\allowbreak{}Token\allowbreak{}Query\allowbreak{}Request]{AuthorizationTokenQueryRequest}}{\mref[Authorization\allowbreak{}Token\allowbreak{}Mgmt\allowbreak{}List\allowbreak{}Response]{AuthorizationTokenMgmtListResponse}}

Returns a page of tokens matching the filter.

\osubsection{revoke-tokens}{\pref{List}$<$\allowbreak{}\pref{Identifier}$>$}{}

Revokes tokens by token reference.

\osubsection{add-encryption-keys}{\mref[Authorization\allowbreak{}Mgmt\allowbreak{}Encryption\allowbreak{}Key\allowbreak{}Registration\allowbreak{}List\allowbreak{}Request]{AuthorizationMgmtEncryptionKeyRegistrationListRequest}}{\mref[Authorization\allowbreak{}Mgmt\allowbreak{}Encryption\allowbreak{}Key\allowbreak{}List\allowbreak{}Response]{AuthorizationMgmtEncryptionKeyListResponse}}

Registers encryption keys on behalf of providers.

\osubsection{remove-encryption-keys}{\pref{List}$<$\allowbreak{}\pref{Name}$>$}{}

Removes the encryption keys of the named providers.

\newpage

\section{Information Model}
\label{sec:model}

Here, all data objects that can be part of the \texttt{authorization\allowbreak{}Token\allowbreak{}Management} service are listed in alphabetic order.
Note that each subsection, which describes one type of object, begins with the \textit{struct} keyword, which is used to denote a collection of named fields, each with its own data type.
As a complement to the explicitly defined types in this section, there is also a list of implicit primitive types in Section \ref{sec:model:primitives}, which are used to represent things like hashes and identifiers.

\msubsection{struct}{AuthorizationMgmtEncryptionKeyListResponse}

List of registered encryption keys.


\noindent\begin{tabularx}{\textwidth}{| >{\raggedright\arraybackslash}p{3.9cm} | >{\raggedright\arraybackslash\small}p{5.3cm} | X |} \hline
\rowcolor{gray!33} Field & Type & Description \\ \hline
entries & \pref{List}$<$\allowbreak{}\mref[Authorization\allowbreak{}Mgmt\allowbreak{}Encryption\allowbreak{}Key\allowbreak{}Response]{AuthorizationMgmtEncryptionKeyResponse}$>$ & Matching records of the current page. \\ \hline
count & \pref{Number} & Total number of matching records (all pages). \\ \hline
\end{tabularx}



\msubsection{struct}{AuthorizationMgmtEncryptionKeyRegistrationListRequest}

Encryption keys to register on behalf of providers.


\noindent\begin{tabularx}{\textwidth}{| >{\raggedright\arraybackslash}p{3.9cm} | >{\raggedright\arraybackslash\small}p{5.3cm} | X |} \hline
\rowcolor{gray!33} Field & Type & Description \\ \hline
list & \pref{List}$<$\allowbreak{}\mref[Authorization\allowbreak{}Mgmt\allowbreak{}Encryption\allowbreak{}Key\allowbreak{}Registration\allowbreak{}Request]{AuthorizationMgmtEncryptionKeyRegistrationRequest}$>$ & The list of request items. \\ \hline
\end{tabularx}



\msubsection{struct}{AuthorizationMgmtEncryptionKeyRegistrationRequest}

Encryption key of one provider.


\noindent\begin{tabularx}{\textwidth}{| >{\raggedright\arraybackslash}p{3.9cm} | >{\raggedright\arraybackslash\small}p{5.3cm} | X |} \hline
\rowcolor{gray!33} Field & Type & Description \\ \hline
system\allowbreak{}Name & \pref{Name} & Provider system the key belongs to. \\ \hline
key & \pref{String} & The encryption key, as a string. \\ \hline
algorithm & \pref{String} & Encryption algorithm of the key, e.g. \texttt{AES/\allowbreak{}ECB/\allowbreak{}PKCS5Padding} or \texttt{AES/\allowbreak{}CBC/\allowbreak{}PKCS5Padding}. \\ \hline
\end{tabularx}



\msubsection{struct}{AuthorizationMgmtEncryptionKeyResponse}

A registered encryption key.


\noindent\begin{tabularx}{\textwidth}{| >{\raggedright\arraybackslash}p{3.9cm} | >{\raggedright\arraybackslash\small}p{5.3cm} | X |} \hline
\rowcolor{gray!33} Field & Type & Description \\ \hline
system\allowbreak{}Name & \pref{Name} & Name of the system, PascalCase (e.g. \texttt{TemperatureSensor}). \\ \hline
raw\allowbreak{}Key & \pref{String} & The registered encryption key. \\ \hline
algorithm & \pref{String} & Encryption algorithm of the key, e.g. \texttt{AES/\allowbreak{}ECB/\allowbreak{}PKCS5Padding} or \texttt{AES/\allowbreak{}CBC/\allowbreak{}PKCS5Padding}. \\ \hline
key\allowbreak{}Additive & \pref{String} & Initialization vector (CBC mode), Base64; empty for ECB. \\ \hline
created\allowbreak{}At & \pref{DateTime} & UTC time of creation. \\ \hline
\end{tabularx}



\msubsection{struct}{AuthorizationTokenGenerationMgmtListRequest}

Tokens to generate on behalf of consumers.


\noindent\begin{tabularx}{\textwidth}{| >{\raggedright\arraybackslash}p{3.9cm} | >{\raggedright\arraybackslash\small}p{5.3cm} | X |} \hline
\rowcolor{gray!33} Field & Type & Description \\ \hline
list & \pref{List}$<$\allowbreak{}\mref[Authorization\allowbreak{}Token\allowbreak{}Generation\allowbreak{}Mgmt\allowbreak{}Request]{AuthorizationTokenGenerationMgmtRequest}$>$ & The list of request items. \\ \hline
\end{tabularx}



\msubsection{struct}{AuthorizationTokenGenerationMgmtRequest}

One token to generate on behalf of a consumer.


\noindent\begin{tabularx}{\textwidth}{| >{\raggedright\arraybackslash}p{3.9cm} | >{\raggedright\arraybackslash\small}p{5.3cm} | X |} \hline
\rowcolor{gray!33} Field & Type & Description \\ \hline
token\allowbreak{}Variant & \pref{String} & Requested token variant, as interface security policy (e.g. \texttt{TIME\_\allowbreak{}LIMITED\_\allowbreak{}TOKEN\_\allowbreak{}AUTH}, \texttt{RSA\_\allowbreak{}SHA512\_\allowbreak{}JSON\_\allowbreak{}WEB\_\allowbreak{}TOKEN\_\allowbreak{}AUTH}). \\ \hline
target\allowbreak{}Type & \pref{String} & Target type: \texttt{SERVICE\_\allowbreak{}DEF} or \texttt{EVENT\_\allowbreak{}TYPE}. \\ \hline
consumer\allowbreak{}Cloud & \pref{Name} & Cloud identifier of the consumer system. \\ \hline
consumer & \pref{Name} & Name of the consumer system. \\ \hline
provider & \pref{Name} & Name of the provider system. \\ \hline
target & \pref{Name} & Target of the policy or token: a service definition or event type. \\ \hline
scope & \pref{Name} & Operation the token or policy is limited to; empty for all operations. \\ \hline
expires\allowbreak{}At & \pref{DateTime} & Requested expiry; default \texttt{token.\allowbreak{}time.\allowbreak{}limit} from now. \\ \hline
usage\allowbreak{}Limit & \pref{Number} & Requested uses; default \texttt{simple.\allowbreak{}token.\allowbreak{}usage.\allowbreak{}limit}. \\ \hline
\end{tabularx}



\msubsection{struct}{AuthorizationTokenMgmtListResponse}

List of generated tokens.


\noindent\begin{tabularx}{\textwidth}{| >{\raggedright\arraybackslash}p{3.9cm} | >{\raggedright\arraybackslash\small}p{5.3cm} | X |} \hline
\rowcolor{gray!33} Field & Type & Description \\ \hline
entries & \pref{List}$<$\allowbreak{}\mref[Authorization\allowbreak{}Token\allowbreak{}Response]{AuthorizationTokenResponse}$>$ & Matching records of the current page. \\ \hline
count & \pref{Number} & Total number of matching records (all pages). \\ \hline
\end{tabularx}



\msubsection{struct}{AuthorizationTokenQueryRequest}

Filter for generated tokens.


\noindent\begin{tabularx}{\textwidth}{| >{\raggedright\arraybackslash}p{3.9cm} | >{\raggedright\arraybackslash\small}p{5.3cm} | X |} \hline
\rowcolor{gray!33} Field & Type & Description \\ \hline
pagination & \mref{Page} & Paging and sorting of the result. \\ \hline
requester & \pref{Name} & Name of the system that requested the token. \\ \hline
token\allowbreak{}Type & \pref{String} & Token type to match. \\ \hline
consumer\allowbreak{}Cloud & \pref{Name} & Cloud identifier of the consumer system. \\ \hline
consumer & \pref{Name} & Name of the consumer system. \\ \hline
provider & \pref{Name} & Name of the provider system. \\ \hline
target\allowbreak{}Type & \pref{String} & Target type: \texttt{SERVICE\_\allowbreak{}DEF} or \texttt{EVENT\_\allowbreak{}TYPE}. \\ \hline
target & \pref{Name} & Target of the policy or token: a service definition or event type. \\ \hline
\end{tabularx}



\msubsection{struct}{AuthorizationTokenResponse}

A generated token with its bookkeeping data.


\noindent\begin{tabularx}{\textwidth}{| >{\raggedright\arraybackslash}p{3.9cm} | >{\raggedright\arraybackslash\small}p{5.3cm} | X |} \hline
\rowcolor{gray!33} Field & Type & Description \\ \hline
token\allowbreak{}Type & \pref{AuthorizationTokenType} & Token type: \texttt{TIME\_\allowbreak{}LIMITED\_\allowbreak{}TOKEN}, \texttt{USAGE\_\allowbreak{}LIMITED\_\allowbreak{}TOKEN}, \texttt{SELF\_\allowbreak{}CONTAINED\_\allowbreak{}TOKEN} or \texttt{TRANSLATION\_\allowbreak{}BRIDGE\_\allowbreak{}TOKEN}. \\ \hline
variant & \pref{String} & Token variant (interface security policy). \\ \hline
token & \pref{Token} & The token, for simple tokens only in hashed form. \\ \hline
token\allowbreak{}Reference & \pref{Identifier} & Reference of the token, used to revoke it. \\ \hline
requester & \pref{Name} & Name of the system that requested the token. \\ \hline
consumer\allowbreak{}Cloud & \pref{Name} & Cloud identifier of the consumer system. \\ \hline
consumer & \pref{Name} & Name of the consumer system. \\ \hline
provider & \pref{Name} & Name of the provider system. \\ \hline
target\allowbreak{}Type & \pref{AuthorizationTargetType} & Target type: \texttt{SERVICE\_\allowbreak{}DEF} or \texttt{EVENT\_\allowbreak{}TYPE}. \\ \hline
target & \pref{Name} & Target of the policy or token: a service definition or event type. \\ \hline
scope & \pref{Name} & Operation the token or policy is limited to; empty for all operations. \\ \hline
created\allowbreak{}At & \pref{DateTime} & UTC time of creation. \\ \hline
usage\allowbreak{}Limit & \pref{Number} & Number of allowed uses (usage-limited tokens). \\ \hline
usage\allowbreak{}Left & \pref{Number} & Remaining uses of a usage-limited token. \\ \hline
expires\allowbreak{}At & \pref{DateTime} & UTC expiry time; no value means no expiry. \\ \hline
\end{tabularx}



\msubsection{struct}{Page}

Paging and sorting; all fields optional.


\noindent\begin{tabularx}{\textwidth}{| >{\raggedright\arraybackslash}p{3.9cm} | >{\raggedright\arraybackslash\small}p{5.3cm} | X |} \hline
\rowcolor{gray!33} Field & Type & Description \\ \hline
page & \pref{Number} & Zero-based page number. \\ \hline
size & \pref{Number} & Number of records per page (at most \texttt{max.\allowbreak{}page.\allowbreak{}size}, default 1000). \\ \hline
direction & \pref{String} & Sort direction, \texttt{ASC} or \texttt{DESC}. \\ \hline
sort\allowbreak{}Field & \pref{String} & Field to sort by. \\ \hline
\end{tabularx}



\subsection{Primitives}
\label{sec:model:primitives}

Types and structures mentioned throughout this document that are assumed to be available to implementations of this service.
The concrete interpretations of each of these types and structures must be provided by any IDD document claiming to implement this service.

\noindent\begin{tabularx}{\textwidth}{| p{4.6cm} | X |} \hline
\rowcolor{gray!33} Type & Description \\ \hline
\pdef{AuthorizationTargetType} & One of \texttt{SERVICE\_\allowbreak{}DEF}, \texttt{EVENT\_\allowbreak{}TYPE}. \\ \hline
\pdef{AuthorizationTokenType} & One of \texttt{TIME\_\allowbreak{}LIMITED\_\allowbreak{}TOKEN}, \texttt{USAGE\_\allowbreak{}LIMITED\_\allowbreak{}TOKEN}, \texttt{SELF\_\allowbreak{}CONTAINED\_\allowbreak{}TOKEN}, \texttt{TRANSLATION\_\allowbreak{}BRIDGE\_\allowbreak{}TOKEN}. \\ \hline
\pdef{DateTime} & Pinpoints a specific moment in time (UTC). \\ \hline
\pdef{Identifier} & A string that identifies a record, e.g. a UUID or a composite identifier. \\ \hline
\pdef{List} & An \textit{array} of a known number of items, each having type A. \\ \hline
\pdef{Name} & A string identifier that is intended to be both human and machine-readable. \\ \hline
\pdef{Number} & An integer number. \\ \hline
\pdef{String} & A chain of characters. \\ \hline
\pdef{Token} & An opaque string used as credential. \\ \hline
\end{tabularx}

\newpage

\bibliographystyle{IEEEtran}
\bibliography{bibliography}

\newpage

\section{Revision History}
\subsection{Amendments}

\noindent\begin{tabularx}{\textwidth}{| p{1cm} | p{3cm} | p{2cm} | X | p{4cm} |} \hline
\rowcolor{gray!33} No. & Date & Version & Subject of Amendments & Author \\ \hline

1 & 2026-10-22 & 5.2.2 & Document created from the Arrowhead LaTeX templates for the 5.2.2 Java Spring implementation, including the local extensions of this workspace. & Jerker Delsing \\ \hline
2 & 2026-11-09 & 5.3.0 & New documentation edition; primary-source references and Go/Rust implementation scope; 5.2.2 retained. & Documentation maintenance \\ \hline
\end{tabularx}

\subsection{Quality Assurance}

\noindent\begin{tabularx}{\textwidth}{| p{1cm} | p{3cm} | p{2cm} | X |} \hline
\rowcolor{gray!33} No. & Date & Version & Approved by \\ \hline

1 & & \arrowversion  &  \\ \hline

\end{tabularx}

\end{document}
